\chapter{Introducere}

\section{Motivație}
Mereu am fost pasionată de călătorii și de organizarea acestora. În era digitală actuală, am observat că există o multitudine de aplicații care oferă informații turistice, care își oferă posibilitatea de a face liste cu destinații sau care marchează destinațiile pe hartă. Insă puține dintre ele se concentrează pe planificarea detaliată a itinerariilor. 

Mi-am dorit să creez o soluție care să ajute utilizatorii să-și organizeze eficient traseele turistice, să aibă posibilitatea de a adăuga rezervări și să ofere un plan zilnic coerent pentru a maximiza experiența de călătorie.

Motivația mea este de a inspira cât mai mulți oameni să descopere diversitatea lumii, sa viziteze cât mai mult fără stresul organizării.

\section{Nevoi pe care aplicația le satisface}
Scopul principal al aplicatiei este de a veni in sprijinul pasionaților de calatorii care nu au timpul necesar sau cărora, pur si simplu, nu le face plăcere procesul de organizare. Aplicația elimină complet efortul de a căuta informații pe site-uri reducând oboseala planificării. Prin preluarea acestei etape utilizatorul se poate axa explusiv pe relaxare si vizitarea țării respective.

Pentru o eficientizare maximă, platforma este concepută să se adapteze stilului fiecărui utilizator. Acesta poate adăuga manual destinații si rezervări pe care le-a făcut deja, sau poate lăsa aplicația să genereze partial sau total itinerariul. Pentru a crea un plan automatizat si personalizat, preferintele sunt colectate într-un mod practic si interactiv: utilizatorul sortează o lista de categorii si le sortează pe trei nivele de prioritate. Astfel traseul generat se va mula pe interesele utilizatorului. 

Nu în ultimul rând, aplicația rezolvă problema centralizării informațiilor. Prin crearea unui cont personalizat utilizatorii scapa de grija rezervărilor, adreselor si notitelor. Totul este adunat intr-un spatiu clar. Din secțiunea principală acesta poate vizualiza si gestiona rapid călătoriile curente sau viitoare (cat si cele trecute). În plus, pentru un ton mai interactiv interfața integrează un calendar în care perioadele de vacanță sunt evidențiate alaturi de o hartă interactiva în care sunt evidențiate țarile explorate.

\section{Aplicații similare}
\ding{117} Google Maps

„Google Maps” este o aplicație destinată căutării de locații și calculării direcțiilor între diferite puncte de interes. Funcționarea acesteia se bazează pe utilizarea unui cont Google asociat, oferind utilizatorului posibilitatea de a căuta destinații atât după denumire și coordonate, cât și prin selectarea dintr-o anumită categorie de interes (restaurante, atracții etc.).

O funcționalitate importantă este posibilitatea de a salva locațiile preferate în liste individuale, facilitând organizarea personalizată. De asemenea, aplicația permite calcularea rutelor între mai multe puncte de oprire. Totuși, un dezavantaj în planificarea itinerariilor complexe este faptul că ordinea locațiilor nu este optimizată automat, aceasta trebuind să fie stabilită manual, uneori în mod aleatoriu, de către utilizator.

\ding{117} Wanderlog

„Wanderlog” este o aplicație de planificare a călătoriilor ce permite utilizatorilor adăugarea de destinații, unități de cazare și atracții turistice, pe care le integrează și le plasează pe o hartă interactivă. Pentru a eficientiza organizarea traseului, aplicația calculează în mod automat distanțele și timpul de deplasare necesar între obiectivele selectate.

Deși oferă funcționalități utile organizării, experiența utilizatorului poate fi afectată de prezența diverselor sfaturi afișate constant sub formă de ferestre de tip pop-up, care pot deveni deranjante pe parcursul utilizării. De asemenea, un aspect limitativ este faptul că aplicația condiționează accesul la foarte multe acțiuni și funcții avansate de achiziționarea unei subscripții plătite.

\ding{117} Roamy

„Roamy” este o aplicație dedicată creării de itinerarii, disponibilă în prezent exclusiv în magazinul de aplicații Apple Store. La inițierea unei noi planificări, utilizatorul trebuie să introducă destinația dorită și numărul de zile alocate călătoriei. Ulterior, traseul poate fi construit fie prin adăugarea manuală a locațiilor, fie generat automat cu ajutorul funcțiilor bazate pe inteligență artificială.

Similar altor aplicații din această categorie, versiunea gratuită prezintă anumite restricții. Pentru a debloca toate funcționalitățile oferite de platformă și pentru a beneficia de o experiență de utilizare continuă, fără întreruperi cauzate de afișarea reclamelor, este necesară plata unui abonament lunar sau anual.


\section{Ce aduce nou}

În ceea ce privește noutățile și avantajele aduse față de alte soluții existente pe piață, aplicația se distinge în primul rând printr-un algoritm de recomandare personalizat. Spre deosebire de alte platforme care se limitează la a sugera cele mai vizitate atracții dintr-un oraș, acest algoritm utilizează un sistem avansat de selecție bazat pe preferințele reale ale utilizatorului. Aceste date sunt colectate într-un mod practic și interactiv: utilizatorul ordonează o listă de categorii de interes, încadrându-le în trei niveluri distincte de prioritate. Mai mult, algoritmul oferă flexibilitatea de a genera itinerariul fie integral, fie parțial, completând inteligent lista de locații deja adăugate manual de către utilizator, în cazul în care acestea nu sunt suficiente pentru a acoperi programul unei zile întregi.

Un alt avantaj major în organizarea logistică este serviciul de adăugare a unității de cazare ca punct oficial de pornire. Această funcționalitate contribuie la o calculare mult mai precisă a traseului zilnic, luând în considerare distanța și timpul exact de deplasare de la cazare până la prima destinație. În completarea acestei funcții, aplicația oferă posibilitatea de a integra rezervări prestabilite direct în itinerariu. Aceste evenimente primesc automat cel mai înalt grad de importanță, iar pentru a asigura respectarea programului, sistemul calculează și adaugă o marjă de eroare temporală - ascunsă în interfața utilizatorului - garantând astfel sosirea la timp la rezervările efectuate.

Pe parcursul desfășurării călătoriei, aplicația permite un management dinamic al traseului prin actualizări în timp real. Utilizatorul are posibilitatea de a bifa locațiile vizitate, iar în cazul în care anumite obiective din program nu au putut fi atinse în zilele anterioare, sistemul oferă opțiunea de a regenera itinerariul, reprogramând inteligent locațiile rămase. 

Nu în ultimul rând, având ca scop principal facilitarea planificării și transformarea călătoriei într-o experiență memorabilă, aplicația integrează elemente specifice de gamificare. Printre acestea se numără teste cu întrebări de cultură generală despre destinații, o hartă interactivă pe care țările vizitate sunt evidențiate cromatic, precum și un calendar vizual unde utilizatorii își pot urmări cu ușurință istoricul călătoriilor.

\section{Structura lucrării}

Această lucrare este structurată în patru capitole, fiecare servind un rol specific în prezentarea temei abordate:

\begin{itemize}
    \item Introducere: Oferă o privire de ansamblu asupra subiectului, motivația mea, nevoi pe care aplicația le satisface, o privire de ansamblu a aplicațiilor similare adăugând ce îmbunătățiri am adus eu
    \item Tehnologii și framework-ul utilizat: este structurat pe subcapitole dedicat fiecărei tehnologii utilizate
    \item Prezentarea aplicației: prezintă întreaga funcționalitate a aplicației, din perspectiva unui utlizator obișnuit. Fiecarui subcapitol îi este asociat o pagina cu feature-urile asociate
    \item Concluzia: sunt reunite concluziile și perspectivele de dezvoltare ale aplicației.
\end{itemize}